\documentclass{article}
\usepackage[margin=2.5cm]{geometry}
\usepackage{amsmath, amsfonts, amssymb}
\usepackage{parskip}
\usepackage{enumitem}
\usepackage{booktabs}
\usepackage{tabularx}
\usepackage{array}
\usepackage[bookmarks,colorlinks,allcolors=blue]{hyperref}

\newcommand{\state}{\boldsymbol{\theta}}
\newcommand{\clamp}{\operatorname{clamp}}

\title{EKF Tracking of ESCNN Weights:\\Unsupervised Soft-Syndrome (\texttt{ekf}) and Supervised (\texttt{ekfi}) Measurements}
\author{}
\date{}

\begin{document}
\maketitle

\section{Unsupervised EKF via soft-syndrome measurements (\texttt{ekf})}
\label{sec:ekf}

\subsection{State-space model}
\label{sec:model}

The ESCNN weights are treated as a latent state $\state_i \in \mathbb{R}^d$ that drifts slowly
from frame to frame, and is tracked using only the channel code's parity structure —
no ground-truth bits are needed.

\subsubsection*{State evolution}

\begin{equation}
\state_i = f(\state_{i-1}) + \mathbf{w}_i, \qquad Q = \operatorname{Cov}(\mathbf{w}_i)
\end{equation}

In the simplest (and most common) case $f$ is the identity, i.e.\ the weights follow a random
walk and the process noise $Q$ controls how fast tracking is allowed to move.

\paragraph{Choice of $f$.} Two natural dynamics models:

\begin{itemize}[itemsep=2pt]
  \item \textbf{Random walk:} $f(\state) = \state$, so $F_i = I$. The weight estimate has no
  pull-back term — under this model, whatever the filter drifts to is held purely by the
  balance of $Q$ and $R$, with no anchor to a known-good solution.
  \item \textbf{First-order Markov / AR(1):}
  $f(\state) = \alpha\,\state + (1-\alpha)\,\state_{\text{pretrained}}$, $0 < \alpha \le 1$,
  so $F_i = \alpha I$. The weights are mean-reverting toward the pretrained solution
  $\state_{\text{pretrained}}$ at rate $1-\alpha$ per frame; $\alpha=1$ recovers the random
  walk exactly. This adds an implicit regularizer that keeps the filter anchored near a
  known-good solution rather than free to wander indefinitely under noisy syndrome
  measurements — useful given that "all checks satisfied" alone does not guarantee the tracked
  weights are still correct (see the earlier discussion on degenerate syndrome solutions).
\end{itemize}

\subsubsection*{Measurement model}

\begin{equation}
\mathbf{y}_i = h_i(\state_i) + \mathbf{v}_i, \qquad R = \operatorname{Cov}(\mathbf{v}_i), \qquad
\mathbf{y}_i = \mathbf{1}_M
\end{equation}

The pseudo-label is the all-ones vector: the filter is driven to make every one of the $M$
parity checks satisfied, which is the only supervision signal available. The measurement
function $h_i$ is the composition

\begin{equation}
h_i \; : \quad \state_i \xrightarrow[\mathbf{X}_i]{\text{ESCNN}} \widetilde{\mathbf{L}}_i
\;\to\; \mathbf{q}_i \;\to\; \mathbf{t}_i \;\to\; \mathbf{p}_i ,
\end{equation}

built from the network's soft output and the parity-check matrix $H$ (bit $n$, check $m$,
entry $H_{mn}$):

\begin{align}
\widetilde{L}_{i,n} &= f_{\state_i}(\mathbf{X}_i) \label{eq:llr} \\[4pt]
q_{i,n} &=
\begin{cases}
\widetilde{L}_{i,n}, & n \text{ transmitted} \\
L_{\max}, & n \text{ filler} \\
0, & n \text{ punctured}
\end{cases} \\[4pt]
t_{i,n} &= \tanh\!\left(\frac{\clamp(q_{i,n}, \pm L_{\max})}{2}\right) \label{eq:t} \\[4pt]
p_{i,m} &= \prod_{n \,:\, H_{mn}=1} t_{i,n}, \qquad
\mathbf{p}_i = [p_{i,1}, \dots, p_{i,M}]^T \label{eq:p}
\end{align}

Filler and punctured positions are fixed constants w.r.t.\ $\state_i$; only the transmitted
bits carry gradient information back to the weights. Noise covariances are taken isotropic,
$Q = \sigma_Q^2 I$, $R = \sigma_R^2 I$.

\subsubsection*{Initialization}

\begin{equation}
\hat{\state}_{0|0} = \state_{\text{pretrained}}, \qquad
P_{0|0} = \sigma_{P_0}^2 I, \qquad
Q = \sigma_Q^2 I, \qquad
R = \sigma_R^2 I
\end{equation}

\begin{itemize}[itemsep=2pt]
  \item $\hat{\state}_{0|0}$ is set to the offline-pretrained ESCNN weights — the filter starts
  from a good point and only tracks drift away from it, rather than learning from scratch.
  \item $P_{0|0}$, $Q$, and $R$ are all initialized as \emph{scaled identity matrices}: a single
  scalar variance times $I$, rather than any structured (e.g.\ per-layer or per-check)
  covariance. This keeps the filter cheap to initialize and gives three interpretable knobs:
  \begin{itemize}
    \item $\sigma_{P_0}^2$ — initial confidence in $\hat{\state}_{0|0}$ (larger $\Rightarrow$
    the first few updates are allowed to move the weights more).
    \item $\sigma_Q^2$ — how fast the weights are allowed to drift per frame (tracking speed
    vs.\ stability trade-off).
    \item $\sigma_R^2$ — how much the syndrome measurement is trusted (larger $\Rightarrow$
    softer, more conservative updates).
    \item $\alpha$ (only if the AR(1) dynamics model is used) — reversion strength toward
    $\state_{\text{pretrained}}$; $\alpha \to 1$ behaves like the random walk, smaller $\alpha$
    anchors the estimate more tightly.
  \end{itemize}
  \item All three scalars are hyperparameters to be tuned (e.g.\ via a validation sweep, or by
  monitoring filter consistency through the normalized innovation
  $\Delta\mathbf{y}_i^T S_i^{-1} \Delta\mathbf{y}_i$, which should track $\chi^2(M)$ if $Q$ and
  $R$ are well-scaled). The scaled-identity choice is a starting point, not a constraint — it
  can later be replaced with a block-diagonal or per-layer form if needed.
\end{itemize}

\subsection{Predict step}
\label{sec:predict}

\begin{align}
\hat{\state}_{i|i-1} &= f\big(\hat{\state}_{i-1|i-1}\big) \label{eq:pred-mean}
  && \text{(propagate the estimate through the dynamics)} \\
F_i &= \left.\frac{\partial f(\state)}{\partial \state}\right|_{\hat{\state}_{i-1|i-1}} \label{eq:pred-F}
  && \Big(\Rightarrow\ \state_i-\hat{\state}_{i|i-1} \approx F_i\big(\state_{i-1}-\hat{\state}_{i-1|i-1}\big)+\mathbf{w}_i\Big) \\
P_{i|i-1} &= F_i \, P_{i-1|i-1} \, F_i^T + Q \label{eq:pred-P}
  && \big(= \operatorname{Cov}(\state_i-\hat{\state}_{i|i-1})\text{: old uncertainty + drift}\big)
\end{align}

Equation~\eqref{eq:pred-F} comes from a first-order Taylor expansion,
$f(\state)\approx f(\hat{\state}_{i-1|i-1}) + F_i(\state-\hat{\state}_{i-1|i-1})$, which makes
the prediction error a linear function of the previous estimation error plus the process
noise. Since $\mathbf{w}_i$ is independent of that error, its covariance is \eqref{eq:pred-P}. For both
dynamics models used here $f$ is affine, so the expansion is exact.

\textbf{What needs to be computed:}
\begin{enumerate}[itemsep=2pt]
  \item Propagate the weight estimate through $f$.
  \begin{itemize}[itemsep=2pt]
    \item Random walk: $\hat{\state}_{i|i-1} = \hat{\state}_{i-1|i-1}$ (no-op).
    \item AR(1): $\hat{\state}_{i|i-1} = \alpha\,\hat{\state}_{i-1|i-1} +
    (1-\alpha)\,\state_{\text{pretrained}}$ — a single scalar-weighted average with the fixed
    pretrained checkpoint.
  \end{itemize}
  \item $F_i$, the Jacobian of the dynamics model. Both proposed choices of $f$ are affine in
  $\state$, so $F_i$ is a known constant, not something that needs autodiff:
  $F_i = I$ for the random walk, $F_i = \alpha I$ for AR(1). The predict covariance is then
  simply $P_{i|i-1} = \alpha^2 P_{i-1|i-1} + Q$ (random walk is the $\alpha=1$ special case).
\end{enumerate}

\subsection{Update step}
\label{sec:update}

\begin{align}
\hat{\mathbf{y}}_{i|i-1} &= h_i\big(\hat{\state}_{i|i-1}\big) = \mathbf{p}_i \label{eq:yhat}
  && \text{(predicted checks)} \\
\Delta \mathbf{y}_i &= \mathbf{1}_M - \mathbf{p}_i \label{eq:innov}
  && \text{(innovation: unmet parity)} \\
H_i &= \left.\frac{\partial h_i(\state)}{\partial \state}\right|_{\hat{\state}_{i|i-1}} \label{eq:H}
  && \Big(\Rightarrow\ \Delta\mathbf{y}_i \approx H_i\big(\state_i-\hat{\state}_{i|i-1}\big)+\mathbf{v}_i\Big) \\
S_i &= H_i P_{i|i-1} H_i^T + R \label{eq:S}
  && \big(= \operatorname{Cov}(\Delta\mathbf{y}_i)\big) \\
K_i &= P_{i|i-1} H_i^T S_i^{-1} \label{eq:K}
  && \big(= \operatorname{Cov}(\state_i,\Delta\mathbf{y}_i)\,\operatorname{Cov}(\Delta\mathbf{y}_i)^{-1}\big) \\
\hat{\state}_{i|i} &= \hat{\state}_{i|i-1} + K_i \, \Delta \mathbf{y}_i \label{eq:upd}
  && \text{(LMMSE correction of } \hat{\state}_{i|i-1}\text{)} \\
P_{i|i} &= P_{i|i-1} - K_i S_i K_i^T \label{eq:Pupd}
  && \text{(uncertainty removed by } \Delta\mathbf{y}_i\text{)}
\end{align}

Equation~\eqref{eq:H} comes from a first-order Taylor expansion,
$h_i(\state)\approx \mathbf{p}_i + H_i(\state-\hat{\state}_{i|i-1})$, which makes the
innovation a linear, noisy measurement of the weight error. Equations~\eqref{eq:S}--\eqref{eq:Pupd} are
then the standard linear-Gaussian update: with
$\operatorname{Cov}(\state_i-\hat{\state}_{i|i-1})=P_{i|i-1}$ and
$\mathbf{v}_i\sim(\mathbf{0},R)$, \eqref{eq:S} and the cross-covariance
$P_{i|i-1}H_i^T$ follow directly from the linearized model.

\textbf{What needs to be computed:}
\begin{enumerate}[itemsep=2pt]
  \item \emph{Forward pass:} run $\state_i \to \widetilde{\mathbf{L}}_i \to \mathbf{q}_i \to
  \mathbf{t}_i \to \mathbf{p}_i$ once to get $\hat{\mathbf{y}}_{i|i-1}$ and the innovation
  $\Delta\mathbf{y}_i$.
  \item \emph{Jacobian $H_i$:} this does \emph{not} need to be hand-derived. It is obtained
  automatically by differentiating the same forward pass with reverse-mode autodiff
  (e.g.\ \texttt{torch.func.jacrev} or \texttt{vjp} over $h_i$ as a function of $\state$).
  Since the number of parity checks $M$ is typically far smaller than the number of weights
  $d$, reverse mode is the efficient direction (cost $\sim M$ backward passes).
  \item \emph{Pruning:} any check row $m$ that includes a punctured bit has $p_{i,m}=0$
  identically, and its corresponding row of $H_i$ is also identically zero. These rows carry
  no information and should be dropped before forming $S_i$. (Refined further in
  Section~\ref{sec:fallback}: in practice almost every row touches a punctured bit, so this
  naive rule alone is not enough — punctured bits are first estimated from the code's own
  structure, and pruning is then based on that estimation's provenance, not simply on whether
  a row touches $\mathcal{P}$.)
  \item \emph{Gain and covariance update:} standard linear algebra once $H_i$ (or its action
  on vectors, if $d$ is too large to form $H_i$ explicitly) is available. For large $d$,
  avoid materializing $P$ densely — use a low-rank or diagonal approximation and matrix-free
  $H_i v$ / $H_i^T v$ products instead.
\end{enumerate}

\subsection{The punctured-bit fallback and its self-referential failure mode}
\label{sec:fallback}

The pruning rule in Section~\ref{sec:update} — drop any check row that touches a punctured bit — is a
simplification that does not survive contact with real operating points. At MCS2,
$\text{num\_res}=96$ ($n_{\text{ldpc}}=4576$, $M=3696$), \emph{every single check row}
touches at least one punctured position, so that rule alone leaves $\mathbf{y}_i$ empty and
the update step vacuous. What is actually implemented is a two-stage scheme: first estimate
the punctured bits from the code's own structure, then prune more carefully than "touches a
punctured bit" — this section derives why the second step is necessary.

\subsubsection*{Erasure-peeling estimate of the punctured bits}

Let $\mathcal{P} \subset \{1,\dots,N\}$ be the punctured bit indices (where $q_{i,n}=0$ by
construction, hence $t_{i,n}=0$). For $n \in \mathcal{P}$ still unresolved ($t_{i,n}=0$), a
check $m$ with $H_{mn}=1$ can send $n$ a message only if $n$ is check $m$'s \emph{sole}
remaining unresolved neighbour:
\begin{equation}
\mu_{m \to n} = 2\operatorname{artanh}\!\Big(\prod_{\substack{n' \ne n \\ H_{mn'}=1}} t_{i,n'}\Big),
\qquad\text{sent iff } \big|\{n' \in \mathcal{P} : H_{mn'}=1,\ t_{i,n'}=0\}\big| = 1 .
\end{equation}
Each punctured bit's estimate sums every message it receives this round,
$L_n = \sum_{m \in \mathcal{C}(n)} \mu_{m\to n}$, $t_{i,n} \leftarrow \tanh(L_n/2)$, where
$\mathcal{C}(n)$ is the (possibly empty) set of checks that contributed to $n$ this round.
Resolving one bit can unlock a check for its neighbours next round, so this is iterated for a
fixed round budget (\texttt{tsyn\_fallback\_iters}), stopping early at a fixpoint.

\subsubsection*{The circularity this introduces}

Define $\text{contributing}(m) = 1$ if check $m$ ever supplied a message to one of its own
neighbours, i.e.\ $m \in \mathcal{C}(n)$ for some $n$. If check $m$ is later scored using that
same $n$'s estimate, it is tested against a value it helped fabricate. Note that this
definition — and the exclusion rule built on it below — is entirely degree-agnostic: \emph{a
check is excluded if it ever contributed a message to resolve one of its own punctured
neighbours, regardless of whether that neighbour's degree is 1 or 22.} Degree only determines
how \emph{likely} a check is to end up excluded (see below), not whether the exclusion
criterion itself applies.

This is not merely inelegant — it is provably vacuous in the worst case. Suppose $n$ has
degree 1 (check $m$ is the \emph{only} row touching it) and gets resolved. Then
$\mathcal{C}(n) = \{m\}$ necessarily, and writing $X := \prod_{n' \ne n,\,H_{mn'}=1} t_{i,n'}$,
\begin{equation}
t_{i,n} = \tanh\!\big(\operatorname{artanh}(X)\big) = X
\quad\Longrightarrow\quad
p_{i,m} = t_{i,n} \cdot X = X^2 \;\ge 0 \ \text{identically.}
\end{equation}
$p_{i,m}$ depends only on $|X|$ — the \emph{magnitude} of check $m$'s real, transmitted
neighbours — and is completely blind to their sign. Flipping the sign of any one real neighbour
(exactly the event a parity check exists to catch) leaves $p_{i,m}$ unchanged. Such a check
rewards confident outputs regardless of correctness; scored into the EKF's innovation, it
biases $\hat{\state}_{i|i}$ toward sharpening the network's confidence rather than correcting
its errors. For $n$ with degree $>1$, resolved from several checks at once, the same leakage is
present but diluted — $t_{i,n}$ mixes $m$'s own contribution with genuinely extrinsic
information from the other resolving checks.

\subsubsection*{Classification of checks: clean / self-referential / dead}

Run the peeling process to its round budget or fixpoint, and classify every check row by the
result:
\begin{itemize}[itemsep=2pt]
  \item \textbf{dead}: at least one neighbour of $m$ is punctured and never resolved
  ($t_{i,n}=0$ for some $n \in \mathcal{P} \cap \{n : H_{mn}=1\}$). Since $t_{i,n}=0$
  identically and does not depend on $\state_i$, $p_{i,m}=0$ identically too — the
  corresponding row of $H_i$ is exactly zero, contributing nothing to the update (harmless,
  but wasted compute if left in).
  \item \textbf{self-referential}: every neighbour of $m$ is resolved, and
  $\text{contributing}(m)=1$ — $m$ helped produce at least one of the values it would be
  tested against.
  \item \textbf{clean}: every neighbour of $m$ is resolved, and $\text{contributing}(m)=0$ — $m$
  only \emph{consumes} estimates that came entirely from other checks.
\end{itemize}
Crucially, a degree-1 punctured bit can never yield a clean check: its only possible
contributor and its only possible consumer are the same row, so if it resolves at all, that
row is self-referential by construction. Degree-1 bits are not the only source of
self-referential rows, though — for a bit resolved from several checks (e.g.\ degree 22 at
MCS2), the checks that did the resolving are self-referential even though the bit's value
itself is legitimately multi-sourced; other checks that merely consume that value, without
having contributed to it, remain clean.

This classification depends only on $H$ and which positions are punctured — never on
$\state_i$ or the frame index $i$, since punctured positions are always exactly $0$ at
initialization regardless of the forward pass. It is therefore computed once, not per update.

\paragraph{Example (MCS2, $\text{num\_res}=96$; $Z=88$, $n_{\text{ldpc}}=4576$).} Three real
rows illustrate all three outcomes:
\begin{itemize}[itemsep=2pt]
  \item \emph{Dead:} bit 3080 (degree 1) is check 2200's only unresolved neighbour candidate,
  but check 2200 also touches bit 72 (punctured, still unresolved that round) — two
  simultaneous unknowns block the update rule, so 3080 never receives a value.
  \item \emph{Self-referential:} bit 4136 (degree 1) is resolved solely by check 3256 (real
  neighbours 959, 1182). With $t_{959}=0.9,\ t_{1182}=-0.7$: $t_{4136}=-0.63$ and
  $p_{3256}=(-0.63)(0.9)(-0.7)=0.3969$. Flipping $t_{1182}\to+0.7$ gives $t_{4136}=0.63$ and
  $p_{3256}=0.3969$ again — unchanged, exactly as predicted by $p=X^2$.
  \item \emph{Clean:} bit 17 (degree 22) is resolved from six independent checks
  ($88,586,982,1376,1604,1885$), none of which is check 371. Check 371 (neighbours
  17, 174, 982, 1251) never contributed to either of its two punctured neighbours, so its
  score genuinely reflects whether four independently-sourced values are jointly consistent.
\end{itemize}
At this operating point: 704/3696 rows clean, 1776/3696 self-referential, 1216/3696 dead.

\subsubsection*{Exclusion rule (supersedes the Section~\ref{sec:update} pruning step)}

Only clean rows are scored:
\begin{equation}
\mathcal{M}_{\text{clean}} = \{\, m : \text{every neighbour of } m \text{ resolved, and }
\text{contributing}(m)=0 \,\}, \qquad
\mathbf{y}_i, \mathbf{p}_i, H_i \ \text{restricted to } \mathcal{M}_{\text{clean}}.
\end{equation}
Dead rows are already harmless (zero row in $H_i$) and self-referential rows are actively
biased, not merely uninformative, so both are dropped — refining "drop rows touching a
punctured bit" in Section~\ref{sec:update} into "drop rows that are dead \emph{or} self-referential." Since
$\mathcal{M}_{\text{clean}}$ is static, this reduces to a fixed row-index mask computed once
from $H$, at negligible runtime cost per update.

\paragraph{Remark: raising the round budget does not help.} Increasing
\texttt{tsyn\_fallback\_iters} beyond 1 round does \emph{not} grow $\mathcal{M}_{\text{clean}}$
— empirically (MCS2, MCS6, MCS12, MCS16; rounds 1 through 20) $|\mathcal{M}_{\text{clean}}|$ is
identical at every round count, while dead rows convert entirely into self-referential ones by
round 2--3. The reason is structural: a row $m$ becomes clean only if \emph{all} of its
punctured neighbours resolve simultaneously, via other checks, in the very round before $m$
would otherwise have become their sole remaining resolver. That either happens on the first
round in which $m$'s unresolved count would drop to zero, or it never happens — once $m$ itself
becomes the sole remaining resolver for one of its neighbours, at any round, $m$ is
self-referential from then on, permanently. More rounds only decide \emph{when} a still-dead
row's fate resolves, never whether it resolves clean.

%======================================================================
\section{Supervised EKF with known calibration bits (\texttt{ekfi})}
\label{sec:ekfi}

\texttt{ekfi} is the supervised counterpart of \texttt{ekf}, and the EKF counterpart of the
SGD genie baseline \texttt{sgdbcei} (the ``i'' stands for \emph{impractical}). It uses the same
tracker (\texttt{EkfParamTracker}), the same state model and the same \texttt{escnn\_ekf\_*}
hyperparameters as Section~\ref{sec:ekf}; only the measurement, and the slots it is taken on,
change.

\subsection{Calibration region}

For every group of $S_{\text{data}}=\texttt{slots\_per\_group}$ scored data slots, the
transmitter sends an additional, separate region of $S_{\text{cal}}=\texttt{calib\_slots\_per\_group}$
calibration slots (\texttt{\_build\_calib\_training\_data} in \texttt{ekf.py}):
\begin{itemize}[itemsep=2pt]
  \item same structure as the data slots: 2 DMRS symbols $+$ 12 payload symbols of fresh,
  CRC+LDPC-coded random bits;
  \item same channel realization $\mathbf{h}$ as the group's data, and the same random draw as
  \texttt{sgdbcei}, so the two modes see identical calibration data;
  \item its own DMRS channel estimate and its own LMMSE prior, which is fed to the ESCNN as the
  augmentation input;
  \item used only for the EKF updates and never scored. The scored data slots are never
  measured in this mode.
\end{itemize}
The calibration payload bits $b_{s,n}\in\{0,1\}$ are assumed known at the receiver; this is what
makes the mode a genie baseline rather than a deployable scheme.

\subsection{Measurement model}

For calibration slot $s$ and user $k$, let $\widetilde{L}_{s,n}(\state)$, $n=1,\dots,N$, be the
ESCNN output LLRs of that slot (convention $\widetilde{L}>0 \Leftrightarrow b=1$), with
$N = 12\, q_m\, B$ payload bits per slot ($B=\texttt{num\_res}$; e.g.\ $N=4608$ for MCS39, 16QAM, at
$B=96$). Each bit gives one measurement, its \emph{agreement} with the known bit:
\begin{align}
a_{s,n}(\state) &= (2b_{s,n}-1)\,\tanh\!\big(\widetilde{L}_{s,n}(\state)/2\big) \in [-1,1]
  && \text{($+1$: confidently correct, $-1$: confidently wrong)} \label{eq:agree}\\
\mathbf{y}_s &= \mathbf{1}_N
  && \text{(target: every bit confidently correct)} \\
\Delta\mathbf{y}_s &= \mathbf{1}_N - \mathbf{a}_s,\quad
\Delta y_{s,n} = 2\,\big|b_{s,n} - \sigma(\widetilde{L}_{s,n})\big|
  && \text{(innovation)} \label{eq:agree-innov}\\
\frac{\partial a_{s,n}}{\partial \state} &= (2b_{s,n}-1)\,\tfrac12\big(1-t_{s,n}^2\big)\,
  \frac{\partial \widetilde{L}_{s,n}}{\partial \state}
  && \text{(row $n$ of $H_s$; $t_{s,n}=\tanh(\widetilde{L}_{s,n}/2)$)} \label{eq:agree-H}
\end{align}
The measurement has exactly the ``target $\mathbf{1}$, innovation $\mathbf{1}-$measurement''
form of \eqref{eq:innov}, so the tracker code is shared unchanged: \eqref{eq:H}--\eqref{eq:Pupd}
apply with $h_i \to \mathbf{a}_s$, $M\to N$, and $R=\sigma_R^2 I_N$. Unlike \eqref{eq:p}, each row
of $H_s$ involves a single network output: because $b_{s,n}$ is known, every bit can be checked
on its own, whereas a parity check only constrains a combination of bits, so each $p_m$ is a
product over all $d_c$ LLRs of the check. The row is still dense in $\state$, since
$\widetilde{L}_{s,n}$ depends on all shared convolution weights. In \eqref{eq:p}, the gradient
reaching bit $n$ is scaled by $\prod_{n'\neq n} t_{n'}$, so one uncertain bit ($t_{n'}\approx0$)
in a check silences that row for all of its bits; in \eqref{eq:agree-H} there is no such
coupling. No clamp, filler or puncturing handling is needed either: the measured bits are
exactly the transmitted ones.

\subsection{Schedule}

Per group: one predict step \eqref{eq:pred-mean}--\eqref{eq:pred-P}, then $S_{\text{cal}}$
sequential update steps, one per calibration slot, with no predict in between. Only then is the
updated network run on the group's scored data slots. In \texttt{ekf}, the group also gets one
predict, but the $S_{\text{data}}$ updates are taken on the scored slots themselves, which are
then scored with the updated weights (test-time adaptation).

\subsection{What one update optimizes}

An EKF update is one Gauss--Newton step, started at the prior mean, on the MAP objective
\begin{equation}
J(\state) = \tfrac{1}{\sigma_R^2}\,\big\|\mathbf{1}-h(\state)\big\|^2
          + \big(\state-\hat{\state}_{i|i-1}\big)^T P_{i|i-1}^{-1}\big(\state-\hat{\state}_{i|i-1}\big).
\end{equation}
For \texttt{ekfi} the data term is
$\tfrac{4}{\sigma_R^2}\sum_n \big(b_n-\sigma(\widetilde{L}_n)\big)^2$, a Brier (squared
probability error) loss, not the BCE used by \texttt{sgdbcei}. One practical consequence: for a
\emph{confidently wrong} bit, $\Delta y_n\approx 2$ but the factor $1-t_n^2$ in
\eqref{eq:agree-H} is $\approx 0$, so the EKF barely moves it, whereas the BCE gradient
$\sigma(\widetilde{L}_n)-b_n$ stays $\approx\pm1$. The same factor appears in every
$\partial p_m/\partial\state$ of \texttt{ekf}.

\subsection{Comparison: \texttt{ekf} vs.\ \texttt{ekfi}}

Numbers are for MCS39 (16QAM, $R=0.21$) with $B=96$ and 12 payload symbols per slot:
$N=12\cdot4\cdot96=4608$ coded bits, $k=967$ information bits $+$ 16 CRC $=983$, which gives
base graph 2 with lifting size $Z=104$ ($k_{\text{ldpc}}=1040$, 57 filler bits,
$n_{\text{ldpc}}=5408$) and $42Z=4368$ parity checks. The clean subset
$\mathcal{M}_{\text{clean}}\subset\{1,\dots,4368\}$ is computed from the parity-check matrix at
start-up and printed as \texttt{fallback check quality: $X$/4368 clean}.

{\small
\renewcommand{\arraystretch}{1.25}
\begin{tabularx}{\textwidth}{|>{\raggedright\arraybackslash}p{3.2cm}
  |>{\raggedright\arraybackslash}X|>{\raggedright\arraybackslash}X|}
\hline
 & \texttt{ekf} (Section~\ref{sec:ekf}) & \texttt{ekfi} (this section) \\
\hline
\multicolumn{3}{|l|}{\textbf{Same in both}} \\
\hline
State, dynamics $f$, $F_i$ & \multicolumn{2}{l|}{ESCNN weights left unfrozen by \texttt{escnn\_load\_freeze}; random walk or AR(1)} \\
\hline
$P_{0|0}$, $Q$, $R$, $\alpha$ & \multicolumn{2}{l|}{Same \texttt{escnn\_ekf\_*} values, same scaled-identity form} \\
\hline
Predict / update equations & \multicolumn{2}{l|}{\eqref{eq:pred-mean}--\eqref{eq:pred-P} and \eqref{eq:H}--\eqref{eq:Pupd}, same tracker code} \\
\hline
Innovation form & \multicolumn{2}{l|}{Target $\mathbf{1}$, innovation $\mathbf{1}-h(\hat{\state})\in[0,2]$ per row} \\
\hline
Predicts per group & \multicolumn{2}{l|}{One} \\
\hline
Jacobian & \multicolumn{2}{l|}{\texttt{torch.func.jacrev}, chunked by \texttt{escnn\_ekf\_jacobian\_chunk\_size}} \\
\hline
\multicolumn{3}{|l|}{\textbf{Different}} \\
\hline
Supervision & None: code parity only & Known transmitted bits (genie) \\
\hline
Slots measured & The scored data slots themselves & A separate calibration region, never scored \\
\hline
Extra transmission & None & $S_{\text{cal}}$ calibration slots per group \\
\hline
Updates per group & $S_{\text{data}}=\texttt{slots\_per\_group}$ & $S_{\text{cal}}=\texttt{calib\_slots\_per\_group}$ \\
\hline
Measurement row & Soft check $p_m=\prod_{n\in\mathcal{N}(m)} t_n$ & Bit agreement $a_n=(2b_n-1)\,t_n$ \\
\hline
Rows per slot (per user) & $|\mathcal{M}_{\text{clean}}|$, a subset of the 4368 checks & $N=12\,q_m B=4608$ \\
\hline
Network outputs per row & $d_c$ (all bits of the check); gradient to each scaled by the others' $t$ & One (the bit itself); no coupling \\
\hline
Sign information & Parity only: flipping an even number of bits in a check leaves $p_m$ unchanged & Direct: each row says which way the bit should go \\
\hline
Puncturing / filler & Erasure peeling + static clean-row mask (Section~\ref{sec:fallback}) & Not applicable \\
\hline
LLR clamp & $\pm 30$ before $\tanh$ & None \\
\hline
Implicit data term & $\sum_m (1-p_m)^2$ & $4\sum_n (b_n-\sigma(\widetilde{L}_n))^2$ (Brier) \\
\hline
Meaning of $\sigma_R^2$ & Noise on a check product & Noise on a bit agreement; same value, different quantity \\
\hline
Size of $S$ & $|\mathcal{M}_{\text{clean}}|\times|\mathcal{M}_{\text{clean}}|$, at most $4368\times4368$ & $4608\times4608$ \\
\hline
ESCNN prior input & \texttt{which\_augment} prior of the data slots & Always the LMMSE prior of the calibration slots \\
\hline
Main failure mode & Converging to confident, parity-consistent but wrong outputs & Limited by $S_{\text{cal}}$; confident errors are hard to undo \\
\hline
\end{tabularx}
}

Two of these differences deserve a note. First, $\sigma_R$ is shared, but it scales measurements
of different kinds and dimensions, so the same value does not give the two modes the same
measurement trust; tuning it separately per mode would make the comparison fairer. Second, the
calibration prior is always LMMSE, even when \texttt{which\_augment} is, e.g.,
\texttt{AUGMENT\_SPHERE}; in that case \texttt{ekfi} adapts the weights on a different prior from
the one used at scoring time.

%======================================================================
\section{The \texttt{cos\_true} diagnostic in \texttt{ekf}}
\label{sec:costrue}

\texttt{ekf} can log, for every update, how well the step it takes agrees with the step that
knowledge of the true bits would have produced. This is a diagnostic only: it never changes the
filter's trajectory.

\subsection{What ``true'' is}

When logging is on (\texttt{log\_train\_every\_epochs} $>0$), \texttt{ekf.py} passes
\texttt{tx\_ref}, the true CRC+LDPC-coded payload bits $b_{s,n}$ of the \emph{same scored data
slots} that \texttt{ekf} is updating on. It is not the calibration region of
Section~\ref{sec:ekfi}: those slots do not exist in \texttt{ekf} mode. The bits are the very ones
the slot is later scored against.

\subsection{How the reference step is computed}

For each data slot $s$, inside the same update call and from the same prior
$(\hat{\state}_{i|i-1}, P_{i|i-1})$ the applied step uses, i.e.\ after the group's predict and
after the syndrome updates of slots $1,\dots,s-1$:
\begin{align}
\mathbf{r}_s(\state) &= \big[(2b_{s,n}-1)\tanh(\widetilde{L}_{s,n}(\state)/2)\big]_{n=1}^{N}
  && \text{(agreement \eqref{eq:agree}, on slot $s$ itself)} \\
H^{\text{ref}}_s &= \left.\frac{\partial \mathbf{r}_s}{\partial\state}\right|_{\hat{\state}_{i|i-1}},
\quad S^{\text{ref}}_s = H^{\text{ref}}_s P_{i|i-1} H^{\text{ref}\,T}_s + R
  && \text{(same $P_{i|i-1}$, same $\sigma_R^2$)} \\
\Delta\state^{\text{ref}}_s &= P_{i|i-1}H^{\text{ref}\,T}_s \big(S^{\text{ref}}_s\big)^{-1}
  \big(\mathbf{1}_N-\mathbf{r}_s(\hat{\state}_{i|i-1})\big)
  && \text{(computed, never applied)} \\
\Delta\state^{\text{syn}}_s &= K_i\,\Delta\mathbf{y}_i
  && \text{(the step actually applied, \eqref{eq:upd})} \\
\texttt{cos\_true} &= \frac{\langle \Delta\state^{\text{syn}}_s, \Delta\state^{\text{ref}}_s\rangle}
  {\|\Delta\state^{\text{syn}}_s\|\,\|\Delta\state^{\text{ref}}_s\|}
  && \text{($+1$ aligned, $0$ unrelated, $<0$ opposed)}
\end{align}
The LLRs are the same ones that form the syndrome measurement: same network input (received
samples plus prior), same real-bit selection, and the same $N=12\,q_m B$ positions, which equal
the LDPC codeword length \texttt{helper.n}. The cost is a second Jacobian with $N$ rows and an
$N\times N$ solve per slot.

\texttt{cos\_true} is therefore a one-step counterfactual: from exactly where the \texttt{ekf}
filter is now, would a genie that knows this slot's bits push the weights the same way? It does
not track an \texttt{ekfi} trajectory, and it compares directions only, not step sizes.

\subsection{Comparison: applied \texttt{ekf} step vs.\ the \texttt{cos\_true} reference}

{\small
\renewcommand{\arraystretch}{1.25}
\begin{tabularx}{\textwidth}{|>{\raggedright\arraybackslash}p{3.2cm}
  |>{\raggedright\arraybackslash}X|>{\raggedright\arraybackslash}X|>{\raggedright\arraybackslash}X|}
\hline
 & \texttt{ekf} applied step & \texttt{cos\_true} reference & \texttt{ekfi} step (Sec.~\ref{sec:ekfi}) \\
\hline
Slots measured & Scored data slot $s$ & Same scored data slot $s$ & Calibration slot, never scored \\
\hline
Ground truth & None & True bits of slot $s$ & True bits of the calibration slot \\
\hline
Measurement & Soft checks $p_m$, clean rows & Bit agreement $r_n$ & Bit agreement $a_n$ \\
\hline
Rows & $|\mathcal{M}_{\text{clean}}|$ & $N$ & $N$ \\
\hline
LLR clamp / puncturing & $\pm30$; peeling + clean mask & None; not needed & None; not needed \\
\hline
Prior $(\hat{\state}, P)$ & \texttt{ekf}'s own & Identical to the applied step & \texttt{ekfi}'s own trajectory \\
\hline
$R$ & $\sigma_R^2 I$ & $\sigma_R^2 I$ (same value) & $\sigma_R^2 I$ (same value) \\
\hline
Applied to weights & Yes & No & Yes \\
\hline
Purpose & Tracking & Diagnostic: is the blind step pointed the right way? & Genie upper-bound baseline \\
\hline
\end{tabularx}
}

\paragraph{Reading the diagnostic.}
\begin{itemize}[itemsep=2pt]
  \item The reference is stronger than \texttt{ekfi}: it knows the bits of the slot being scored,
  whereas \texttt{ekfi} only knows calibration bits from the same channel.
  \item The cosine is Euclidean over all tracked weights, so the parameter groups with the
  largest step components (e.g.\ the element-wise scales vs.\ the convolution taps) dominate it.
  A per-group cosine, computed separately for the scale vector and for each convolution layer,
  would show which part of the network the blind step gets right.
  \item Logging the norm ratio $\|\Delta\state^{\text{syn}}\|/\|\Delta\state^{\text{ref}}\|$
  alongside the cosine would separate ``right direction, too small'' from ``wrong direction''.
\end{itemize}

\end{document}
